% =============================================================
% commands.tex — ユーティリティコマンド
% =============================================================

% ── 汎用バッジ ──
\NewDocumentCommand{\badge}{m m}{%
  \tikz[baseline=(b.base)]{
    \node[fill=#1, text=white, inner xsep=6pt, inner ysep=2pt,
          font=\sffamily\scriptsize\bfseries, rounded corners=1pt] (b) {#2};
  }%
}

% ── 難易度バッジ ──
\NewDocumentCommand{\diffbadge}{m}{%
  \ifstrequal{#1}{基}{\badge{BadgeKiso}{基\,基礎}}{%
  \ifstrequal{#1}{標}{\badge{BadgeHyojun}{標\,標準}}{%
  \ifstrequal{#1}{強}{\badge{BadgeKyoka}{強\,強化}}{%
  \ifstrequal{#1}{発}{\badge{BadgeHatten}{発\,発展}}{%
    \badge{BadgeHyojun}{#1}%
  }}}}%
}

% ── 出典バッジ ──
\NewDocumentCommand{\sourcebadge}{m}{%
  \ifstrequal{#1}{original}{\badge{OrigBadge}{オリジナル}}{%
    \badge{OfficialBadge}{#1}%
  }%
}

% ── 強調タグ ──
\newcommand{\termtag}[1]{\textcolor{Teal}{\textbf{#1}}}

% ── 比較表テンプレ ──
\NewDocumentCommand{\comptable}{m m}{%
  \begin{center}
  \small
  \renewcommand{\arraystretch}{1.2}%
  \arrayrulecolor{HairLine}%
  \begin{tabularx}{0.96\linewidth}{@{}>{\color{SubBlue}\bfseries\sffamily\footnotesize}l X X@{}}
    \arrayrulecolor{HikakuBar}\toprule
    #1 \\
    \arrayrulecolor{HairLine}\midrule
    #2 \\
    \arrayrulecolor{HikakuBar}\bottomrule
  \end{tabularx}
  \arrayrulecolor{black}%
  \end{center}%
}

% ── PART 区切りページ ──
\NewDocumentCommand{\partpage}{m m}{%
  \clearpage
  \thispagestyle{empty}%
  \null\vfill
  \begin{center}
    {\color{Navy!15}\rule{0.6\linewidth}{0.4pt}}\par
    \vspace{16mm}%
    {\sffamily\Large\color{Muted}PART\enspace #1}\par
    \vspace{10mm}%
    {\fontsize{28pt}{34pt}\selectfont\bfseries\color{Navy}#2\par}%
    \vspace{10mm}%
    {\color{AccentBar}\rule{0.2\linewidth}{2.5pt}}\par
  \end{center}
  \vfill\null
  \clearpage
}

% ── 章扉メタ情報 ──
% 難易度の星
\newcommand{\diffstars}[1]{%
  \foreach \i in {1,...,4}{%
    \ifnum\i>#1\relax
      {\color{HairLine}$\star$}%
    \else
      {\color{Teal}$\bigstar$}%
    \fi
  }%
}

% #1=スキルリスト, #2=目安時間, #3=難易度(1-4), #4=前提知識
\NewDocumentCommand{\chaptermeta}{m m m m}{%
  \begin{tcolorbox}[
    enhanced, breakable,
    colback=LeadBG, colframe=HairLine,
    boxrule=0.4pt, arc=2pt,
    left=14pt, right=14pt, top=12pt, bottom=12pt,
    before skip=4pt, after skip=16pt,
  ]
    {\sffamily\small\bfseries\color{Teal}この章で身につく力}\par\vspace{4pt}%
    {\small\color{Ink}%
      \begin{itemize}[leftmargin=1.2em, itemsep=1pt, topsep=2pt]
        #1
      \end{itemize}
    }%
    \vspace{6pt}%
    {\color{HairLine}\rule{\linewidth}{0.3pt}}\par\vspace{6pt}%
    \noindent
    \begin{minipage}[t]{0.30\linewidth}
      {\sffamily\scriptsize\color{Muted}目安時間}\par
      {\large\bfseries\color{Navy}#2}
    \end{minipage}%
    \begin{minipage}[t]{0.30\linewidth}
      {\sffamily\scriptsize\color{Muted}難易度}\par
      \diffstars{#3}
    \end{minipage}%
    \begin{minipage}[t]{0.40\linewidth}
      {\sffamily\scriptsize\color{Muted}前提知識}\par
      {\footnotesize\color{Ink}#4}
    \end{minipage}
  \end{tcolorbox}
}

% ── 次章への橋渡し ──
\newcommand{\nextchapter}[1]{%
  \par\addvspace{12pt}%
  {\color{AccentBar}\rule{0.3\linewidth}{2pt}}\par\vspace{4pt}%
  {\small\bfseries\color{Navy}次の章では：}%
  {\small\color{Ink}#1\par}%
  \vspace{8pt}%
}

% ── セクション区切りバー ──
\newcommand{\secbar}{%
  \par\vspace{8pt}%
  {\color{TopBar}\rule{0.38\linewidth}{4pt}}%
  {\color{AccentBar}\rule{0.18\linewidth}{4pt}}\par
  \vspace{8pt}%
}

% ── 章末演習見出し ──
\newcommand{\shoumatsuhead}[1]{%
  \par\addvspace{14pt}%
  \noindent\badge{Navy}{#1}%
  \hspace{6pt}%
  {\color{HairLine}\leaders\hrule height 0.4pt\hfill\null}\par
  \vspace{6pt}%
}
